\documentclass{article}
\usepackage[utf8]{inputenc}
\usepackage{geometry}
\usepackage{datetime2}
\usepackage{xparse}
\usepackage{amsmath}
\usepackage{booktabs}
\usepackage{array}
\usepackage{hyperref}
\usepackage{newunicodechar}
\usepackage{float}
\newunicodechar{μ}{\textmu}


 \geometry{
 a4paper,
 total={170mm,257mm},
 left=20mm,
 top=20mm,
 }

 
 \usepackage{graphicx}
 \usepackage{titling}

 \title{Assignment 3: Fuzzy Logic Function\\ Perancangan Fuzzy Inference System (FIS) untuk Penentuan Harga Paket PV
}
% \author{Muhammad Qomaruz Zaman}
\date{11 March 2025}
 
 \usepackage{fancyhdr}
\fancypagestyle{plain}{%  the preset of fancyhdr 
    \fancyhf{} % clear all header and footer fields
    \fancyfoot[R]{}
    \fancyfoot[L]{}
    \fancyhead[L]{\includegraphics[width=2cm]{ITS_pojok.pdf}}
    \fancyhead[R]{ \DTMnow}
}
\makeatletter
\def\@maketitle{%
  \newpage
  \null
  \vskip 1em%
  \begin{center}%
  \let \footnote \thanks
    {\LARGE \@title \par}%
    \vskip 1em%
    %{\large \@date}%
  \end{center}%
  \par
  \vskip 1em}
\makeatother

\usepackage{lipsum}  
% \usepackage{cmbright} % Change the font by zee

\begin{document}

\maketitle


\noindent\begin{tabular}{@{}ll}
     Student        : & Een Hutama Putra\\
     NRP            : & 6022251023\\
     Teacher        : &  Muhammad Qomaruz Zaman, S.T., M.T., Ph.D.\\
     Class name     : & Artificial Intelegent class (X)\\
     %YouTube video link: & \url{https://youtu.be/XnGeqjbIusw?si=C58i2F8sJz37gfQx} \\
     GitHub link    : & \url{https://its.id/m/Task3AIeen}\\
     Overleaf link  : & \url{https://www.overleaf.com/project/69b9cc5d7f56ecd3715dea6c}\\
     Youtube Link    : & \url{https://youtu.be/gPGAudOZnTU}
\end{tabular}\\
    
\vspace{0,5cm}


\section{Pendahuluan}
Dokumen ini menjelaskan perancangan \textit{Fuzzy Inference System (FIS)} menggunakan model Sugeno orde-0 untuk menentukan harga paket sistem Photovoltaic (PV).
\begin{itemize}
    \item Luasan area PV
    \item Daya yang dihasilkan
\end{itemize}

\section{Struktur Sistem}
\begin{itemize}
    \item Tipe FIS: Sugeno orde-0
    \item Metode inferensi: AND (perkalian)
    \item Output: Singleton (konstanta)
    \item Defuzzifikasi: Weighted Average
\end{itemize}

\section{Variabel Input}

\subsection{Luasan Area PV (\texorpdfstring{$A$}{A})}

\begin{center}
\begin{tabular}{|c|c|c|}
\hline
Linguistik & Tipe MF & Parameter \\
\hline
Sedikit & Trapesium & (0, 0, 200, 400) \\
Sedang & Segitiga & (300, 500, 700) \\
Banyak & Trapesium & (600, 800, 1000, 1000) \\
\hline
\end{tabular}
\end{center}

Rentang: 0 -- 1000 m$^2$
\vspace{0,5cm}
\begin{figure}[H]
    \centering
    \includegraphics[width=0.5\linewidth]{Area PV.png}
    \caption{Membership fuzzy logic Area PV.}
    \label{fig:enter-label}
\end{figure}

\subsection{Daya PV ($P$)}

\begin{center}
\begin{tabular}{|c|c|c|}
\hline
Linguistik & Tipe MF & Parameter \\
\hline
Kecil & Trapesium & (0, 0, 100, 200) \\
Menengah & Segitiga & (150, 250, 350) \\
Besar & Trapesium & (300, 400, 500, 500) \\
\hline
\end{tabular}
\end{center}

Rentang: 0 -- 500 kW
\vspace{0,5cm}
\begin{figure}[H]
    \centering
    \includegraphics[width=0.5\linewidth]{Daya PV.png}
    \caption{Membership fuzzy logic Daya PV.}
    \label{fig:enter-label}
\end{figure}
\section{Variabel Output}

\begin{center}
\begin{tabular}{|c|c|c|}
\hline
Paket & Tipe Output & Nilai Singleton \\
\hline
Basic & Konstanta & 1 \\
Premium & Konstanta & 2 \\
Platinum & Konstanta & 3 \\
\hline
\end{tabular}
\end{center}
\vspace{0,5cm}
\begin{figure}[H]
    \centering
    \includegraphics[width=0.5\linewidth]{Output member.png}
    \caption{Output.}
    \label{fig:enter-label}
\end{figure}
\section{Aturan Fuzzy (Rule Base)}
\begin{enumerate}
    \item IF Area sedikit AND Daya kecil THEN Basic
    \item IF Area sedikit AND Daya menengah THEN Basic
    \item IF Area sedikit AND Daya besar THEN Premium
    \item IF Area sedang AND Daya kecil THEN Basic
    \item IF Area sedang AND Daya menengah THEN Premium
    \item IF Area sedang AND Daya besar THEN Platinum
    \item IF Area banyak AND Daya kecil THEN Premium
    \item IF Area banyak AND Daya menengah THEN Platinum
    \item IF Area banyak AND Daya besar THEN Platinum
\end{enumerate}
\vspace{0,5cm}
\begin{center}
\begin{tabular}{|c|c|c|c|}
\hline
No & Area & Daya & Output \\
\hline
1 & Sedikit & Kecil & Basic \\
2 & Sedikit & Menengah & Basic \\
3 & Sedikit & Besar & Premium \\
4 & Sedang & Kecil & Basic \\
5 & Sedang & Menengah & Premium \\
6 & Sedang & Besar & Platinum \\
7 & Banyak & Kecil & Premium \\
8 & Banyak & Menengah & Platinum \\
9 & Banyak & Besar & Platinum \\
\hline
\end{tabular}
\end{center}

\section{Inferensi Sugeno}

\begin{equation}
w_i = \mu_{Area}(A) \times \mu_{Daya}(P)
\end{equation}

\begin{equation}
z_i = \text{konstanta (1, 2, atau 3)}
\end{equation}

\section{Defuzzifikasi}

\begin{equation}
z = \frac{\sum_{i=1}^{n} w_i z_i}{\sum_{i=1}^{n} w_i}
\end{equation}

\section{Contoh Perhitungan}

Misal:
\begin{itemize}
    \item Area = 450 m$^2$
    \item Daya = 320 kW
\end{itemize}

Derajat keanggotaan:
\begin{itemize}
    \item $\mu_{Sedikit} = 0.2$, $\mu_{Sedang} = 0.7$
    \item $\mu_{Menengah} = 0.4$, $\mu_{Besar} = 0.6$
\end{itemize}

Bobot aturan:
\begin{align*}
w_1 &= 0.7 \times 0.4 = 0.28 \\
w_2 &= 0.7 \times 0.6 = 0.42 \\
w_3 &= 0.2 \times 0.4 = 0.08 \\
w_4 &= 0.2 \times 0.6 = 0.12
\end{align*}

Perhitungan output:
\begin{equation}
z = \frac{(0.28 \times 2) + (0.42 \times 3) + (0.08 \times 1) + (0.12 \times 2)}{0.28 + 0.42 + 0.08 + 0.12}
\end{equation}

\begin{equation}
z = \frac{2.14}{0.9} \approx 2.38
\end{equation}
\vspace{0,5cm}
\begin{figure}[H]
    \centering
    \includegraphics[width=0.5\linewidth]{Hasil Output.png}
    \caption{Output.}
    \label{fig:enter-label}
\end{figure}

\section{Interpretasi Output}

\begin{center}
\begin{tabular}{|c|c|}
\hline
Nilai $z$ & Paket \\
\hline
1.0 -- 1.5 & Basic \\
1.5 -- 2.5 & Premium \\
2.5 -- 3.0 & Platinum \\
\hline
\end{tabular}
\end{center}

Hasil menunjukkan bahwa sistem memilih paket \textbf{Premium}.

\section{Kesimpulan}
FIS Sugeno orde-0 memberikan metode yang sederhana dan efisien untuk menentukan harga paket PV berdasarkan parameter teknis. Model ini dapat dikembangkan lebih lanjut dengan menambahkan variabel input lain seperti irradiance atau biaya investasi.




\end{document}